% MATH 100 official midterm examination.
% All questions in this version are distinct from the Lecture 8 mock exam.
% Student version: compile this file directly.
% Instructor version: compile solutions.tex, which defines \instructorversion.
\documentclass[11pt]{article}

\usepackage[a4paper,margin=18mm,headheight=15pt]{geometry}
\usepackage{amsmath,amssymb,mathtools}
\usepackage{array,booktabs}
\usepackage{enumitem}
\usepackage{needspace}
\usepackage{xcolor}
\usepackage{fancyhdr}
\usepackage{lastpage}

\definecolor{solutionblue}{RGB}{0,76,170}
\newif\ifshowsolutions
\ifdefined\instructorversion
  \showsolutionstrue
\else
  \showsolutionsfalse
\fi

\newcommand{\R}{\mathbb{R}}
\newcommand{\vect}[1]{\mathbf{#1}}
\newcommand{\Col}{\operatorname{Col}}
\newcommand{\rank}{\operatorname{rank}}
\newcommand{\pts}[1]{\hfill\textbf{[#1 pts]}}
\newcommand{\answerline}{\ifshowsolutions\else\par\vspace{0.48in}\hrulefill\par\fi}
\newcommand{\workarea}[1]{\ifshowsolutions\else\vspace{#1}\fi}
\newcommand{\mcneedspace}{\ifshowsolutions\Needspace{12\baselineskip}\fi}
\newcommand{\solution}[1]{%
  \ifshowsolutions
    \par\medskip
    {\color{solutionblue}\textbf{Solution.} #1}
    \par\medskip
  \fi}

\def\CourseTitle{Linear Algebra}
\def\ExamName{Midterm Examination}
\def\CourseNumber{MATH 100}
\def\ExamDate{October 28}
\def\Term{Fall 2026}
\def\ExamDuration{2 hours}
\def\TotalQuestions{21}
\def\TotalPoints{100}

\pagestyle{fancy}
\fancyhf{}
\lhead{MATH 100: Linear Algebra}
\rhead{Midterm Examination}
\cfoot{\thepage}

\fancypagestyle{coverpage}{%
  \fancyhf{}
  \cfoot{Page \thepage\ of \pageref{LastPage}}
  \rfoot{Please go on to the next page\ldots}
  \renewcommand{\headrulewidth}{0pt}
  \renewcommand{\footrulewidth}{0pt}}

\setlist[itemize]{leftmargin=1.5em,itemsep=0.25em,topsep=0.35em}
\setlist[enumerate,1]{leftmargin=*,itemsep=0.9em}
\setlist[enumerate,2]{leftmargin=2.0em,itemsep=0.2em,topsep=0.3em}

\begin{document}

\thispagestyle{coverpage}

\begin{center}
  \Large
  {{\CourseNumber} - {\CourseTitle}}\\
  \vspace{5pt}
  \textbf{\ExamName}\\
  \vspace{5pt}
  Term: \enspace {\Term} \qquad Date: \enspace {\ExamDate}
  \ifshowsolutions
    \\\vspace{8pt}
    {\large\color{solutionblue}\bfseries Instructor solution version}
  \fi
\end{center}

\vspace{10pt}
\hrule
\vspace{10pt}

\begin{center}
\vspace{10pt}
\textbf{Please read the following instructions carefully.}
\vspace{10pt}
\begin{itemize}
  \item Write your name, student ID, and sign the cover page.
  \item The duration of this examination is \textbf{\ExamDuration}.
  \item This is a closed-book examination. No phones, notes, or other electronic devices are allowed.
  \item This question booklet contains \textbf{\TotalQuestions} questions, \textbf{\pageref{LastPage}} pages (including the cover), for a total of \textbf{\TotalPoints} points. Check that no pages are missing.
  \item Show all of your work. Marks are awarded for the work, not just the final answer.
  \item Read the instructions for individual questions carefully and clearly identify each answer.
  \item If you need more space, clearly label work completed on additional paper.
  \item Return all examination and draft papers upon completion.
  \item \textbf{\large Good Luck!}
\end{itemize}
\end{center}

\vspace{2cm}
\noindent
\begin{minipage}{0.58\linewidth}
  \raggedright
  \makebox[3.6in]{\textbf{\large Student Name:} \hrulefill} \\[26pt]
  \makebox[3.6in]{\textbf{\large Student ID:} \hrulefill} \\[26pt]
  \makebox[3.6in]{\textbf{\large Signature:} \hrulefill} \\[45pt]
\end{minipage}

\newpage

\section*{Section I. Multiple choice \hfill 20 points}

\begingroup\small
Choose the best answer. Each question is worth 2 points.

\begin{enumerate}
\mcneedspace
\item Which matrix is in row echelon form?
\begin{center}
\begin{tabular}{@{}p{.45\linewidth}p{.45\linewidth}@{}}
A. $\begin{bmatrix}1&0\\0&0\\0&1\end{bmatrix}$ &
B. $\begin{bmatrix}1&2&0\\0&0&1\\0&0&0\end{bmatrix}$ \\[1.5em]
C. $\begin{bmatrix}0&1\\1&0\end{bmatrix}$ &
D. $\begin{bmatrix}1&0\\1&1\end{bmatrix}$
\end{tabular}
\end{center}
\solution{\textbf{B.} The leading entries move to the right as the rows descend, and all entries below each pivot are zero.}

\mcneedspace
\item Which row operation changes
\[
\begin{bmatrix}1&2\\3&7\end{bmatrix}
\quad\text{into}\quad
\begin{bmatrix}1&2\\0&1\end{bmatrix}?
\]
\begin{enumerate}[label=\Alph*.]
  \item $R_1\leftarrow R_1-3R_2$
  \item $R_2\leftarrow R_2-3R_1$
  \item $R_2\leftarrow R_2+3R_1$
  \item $R_1\leftrightarrow R_2$
\end{enumerate}
\solution{\textbf{B.} The new second row is $(3,7)-3(1,2)=(0,1)$.}

\mcneedspace
\item If $T(x,y)=(2x-y,\ x+3y)$, then $T(1,-2)$ equals
\begin{enumerate}[label=\Alph*.]
  \item $(0,-5)$
  \item $(4,-5)$
  \item $(4,7)$
  \item $(-4,-5)$
\end{enumerate}
\solution{\textbf{B.} Substitution gives $T(1,-2)=(2+2,1-6)=(4,-5)$.}

\mcneedspace
\item Suppose Matrix $A$ is $2\times3$, $B$ is $3\times2$, and $C$ is $2\times2$. Which expression is \emph{not} defined?
\begin{enumerate}[label=\Alph*.]
  \item $A+B$
  \item $AB$
  \item $BA$
  \item $CA$
\end{enumerate}
\solution{\textbf{A.} Matrix addition requires equal dimensions, but $A$ and $B$ have different sizes.}

\mcneedspace
\item If $A$ is invertible and $XA=C$, then
\begin{enumerate}[label=\Alph*.]
  \item $X=A^{-1}C$
  \item $X=CA^{-1}$
  \item $X=AC^{-1}$
  \item $X=C^{-1}A$
\end{enumerate}
\solution{\textbf{B.} Multiply on the right by $A^{-1}$: $XAA^{-1}=CA^{-1}$.}

\mcneedspace
\item The determinant of the matrix
\[
\begin{bmatrix}
1&2&1\\
2&1&3\\
1&1&2
\end{bmatrix}
\]
is
\begin{enumerate}[label=\Alph*.]
  \item $-2$
  \item $2$
  \item $-4$
  \item $4$
\end{enumerate}
\solution{\textbf{A.} Expanding along the first row gives
$1(2-3)-2(4-3)+(2-1)=-1-2+1=-2$.}

\mcneedspace
\item If $A$ is $3\times3$ and $\det A=-3$, then $\det(2A^T)$ equals
\begin{enumerate}[label=\Alph*.]
  \item $-6$
  \item $-12$
  \item $-18$
  \item $-24$
\end{enumerate}
\solution{\textbf{D.} $\det(2A^T)=2^3\det(A^T)=8\det A=-24$.}

\mcneedspace
\item Let
\[
\vect v_1=\begin{bmatrix}1\\0\\2\end{bmatrix},
\qquad
\vect v_2=\begin{bmatrix}0\\1\\-1\end{bmatrix}.
\]
Which vector belongs to $\operatorname{Span}(\vect v_1,\vect v_2)$?
\begin{center}
\begin{tabular}{@{}p{.45\linewidth}p{.45\linewidth}@{}}
A. $\begin{bmatrix}2\\-1\\5\end{bmatrix}$ &
B. $\begin{bmatrix}1\\1\\0\end{bmatrix}$ \\[1.5em]
C. $\begin{bmatrix}0\\2\\1\end{bmatrix}$ &
D. $\begin{bmatrix}3\\0\\5\end{bmatrix}$
\end{tabular}
\end{center}
\solution{\textbf{A.} The vector equals $2\vect v_1-\vect v_2$.}

\mcneedspace
\item A matrix has $7$ columns, and its reduced row echelon form has $2$ pivots. How many free variables are there?
\begin{enumerate}[label=\Alph*.]
  \item $2$
  \item $5$
  \item $7$
  \item $9$
\end{enumerate}
\solution{\textbf{B.} There are $7-2=5$ nonpivot columns, so there are five free variables.}

\mcneedspace
\item Assuming all blocks are conformable, the lower-left block of
\[
\begin{bmatrix}P&0\\Q&R\end{bmatrix}
\begin{bmatrix}S&0\\U&V\end{bmatrix}
\]
is
\begin{enumerate}[label=\Alph*.]
  \item $PS$
  \item $PV$
  \item $QU+RV$
  \item $QS+RU$
\end{enumerate}
\solution{\textbf{D.} The lower-left block is obtained from the second block row and first block column.}
\end{enumerate}
\endgroup

\ifshowsolutions\else\newpage\fi
\section*{Section II. Brief answers \hfill 20 points}

Give a short answer and enough justification to support it. Each question is worth 4 points.

\begin{enumerate}
\item For the matrix
\[
E=\begin{bmatrix}
1&0&2&0&-1\\
0&1&-3&0&4\\
0&0&0&1&2\\
0&0&0&0&0
\end{bmatrix},
\]
state $\rank E$ and identify the free-variable columns.
\answerline
\solution{The pivot columns are $1,2,4$, so $\boxed{\rank E=3}$ and the free-variable columns are $\boxed{3\text{ and }5}$.}

\item Let $T:\R^2\to\R^2$ be defined by
\[
T(x,y)=(x+2y,\ 3x-y).
\]
What is the matrix of the transformation $T$?
\answerline
\solution{The columns are $T(1,0)$ and $T(0,1)$, so the standard matrix is $\boxed{\begin{bmatrix}1&2\\3&-1\end{bmatrix}}$.}

\item Suppose $A$ is invertible and
\[
A(X-I)=B.
\]
Use $A$ and $B$ to represent $X$.
\answerline
\solution{Multiply on the left by $A^{-1}$: $X-I=A^{-1}B$. Hence $\boxed{X=I+A^{-1}B}$.}

\item Suppose $A$ and $B$ are invertible $3\times3$ matrices with $\det A=2$ and $\det B=-3$. Find $\det(A^{-1}B)$.
\answerline
\solution{$\det(A^{-1}B)=\det(A^{-1})\det(B)=\frac12(-3)=\boxed{-\frac32}$.}

\item If $A$ is $4\times3$ and $B$ is $3\times2$, what is the size of $AB$?
\answerline
\solution{The inner dimensions agree. The outer dimensions give $\boxed{4\times2}$.}
\end{enumerate}

\newpage
\section*{Section III. Exercises \hfill 60 points}

Show the main steps. Point values are shown for each exercise.

\begin{enumerate}
\item Solve the system by Gaussian elimination. Write the final answer as an ordered triple. \pts{10}
\[
\begin{cases}
x+y+z=4,\\
2x-y+z=8,\\
-x+2y+z=-1.
\end{cases}
\]
\workarea{2.2in}
\solution{Starting from the augmented matrix, use $R_2\leftarrow R_2-2R_1$ and $R_3\leftarrow R_3+R_1$:
\[
\left[\begin{array}{ccc|c}
1&1&1&4\\2&-1&1&8\\-1&2&1&-1
\end{array}\right]
\longrightarrow
\left[\begin{array}{ccc|c}
1&1&1&4\\0&-3&-1&0\\0&3&2&3
\end{array}\right].
\]
Adding the second row to the third gives $z=3$. Then $-3y-z=0$ gives $y=-1$, and the first equation gives $x=2$. Thus
\[
\boxed{(x,y,z)=(2,-1,3)}.
\]}

% \item Compute $AB$. \pts{10}
% \[
% A=\begin{bmatrix}1&-2\\0&3\\4&1\end{bmatrix},
% \qquad
% B=\begin{bmatrix}2&1&0\\-1&2&5\end{bmatrix}.
% \]
% \workarea{2.0in}
% \solution{
% \[
% AB=
% \begin{bmatrix}
% 1(2)+(-2)(-1)&1(1)+(-2)(2)&1(0)+(-2)(5)\\
% 0(2)+3(-1)&0(1)+3(2)&0(0)+3(5)\\
% 4(2)+1(-1)&4(1)+1(2)&4(0)+1(5)
% \end{bmatrix}
% =\boxed{\begin{bmatrix}4&-3&-10\\-3&6&15\\7&6&5\end{bmatrix}}.
% \]}


\item Let
\[
A=\begin{bmatrix}0&1\\-2&3\end{bmatrix},
\qquad
B=\begin{bmatrix}1&-1\\4&0\end{bmatrix}.
\]
Compute $A^2+AB+I$. \pts{10}
\workarea{2.0in}
\solution{First,
\[
A^2=\begin{bmatrix}-2&3\\-6&7\end{bmatrix},
\qquad
AB=\begin{bmatrix}4&0\\10&2\end{bmatrix}.
\]
Therefore
\[
A^2+AB+I_2
=\begin{bmatrix}-2&3\\-6&7\end{bmatrix}
+\begin{bmatrix}4&0\\10&2\end{bmatrix}
+\begin{bmatrix}1&0\\0&1\end{bmatrix}
=\boxed{\begin{bmatrix}3&3\\4&10\end{bmatrix}}.
\]}
\newpage
% \item Define $T:\R^3\to\R^2$ by
% \[
% T(x,y,z)=(x-y+2z,\ 2x+y-z).
% \]
% \begin{enumerate}[label=(\alph*)]
%   \item Find the standard matrix of $T$. \hfill[2 pts]
%   \item Compute $T(1,2,-1)$. \hfill[2 pts]
%   \item Determine whether $\vect b=(3,1)$ belongs to the range of $T$. If it does, find one vector $\vect x$ such that $T(\vect x)=\vect b$. \hfill[3 pts]
%   \item Determine whether $T$ is one-to-one and whether it is onto. \hfill[3 pts]
% \end{enumerate}
% \workarea{2.1in}
% \solution{
% \begin{enumerate}[label=(\alph*)]
%   \item $\displaystyle [T]=\boxed{\begin{bmatrix}1&-1&2\\2&1&-1\end{bmatrix}}$.
%   \item $\displaystyle T(1,2,-1)=\boxed{(-3,5)}$.
%   \item Yes. For example, $\vect x=(1,0,1)$ gives $T(1,0,1)=(3,1)$, so $\boxed{\vect b\text{ belongs to the range of }T}$.
%   \item The standard matrix has a pivot in each row, so $T$ is onto $\R^2$. It has only two pivots but three columns, so it is not one-to-one.
% \end{enumerate}}

% \newpage
\item Find an $LU$ factorization $A=LU$, where $L$ is unit lower triangular. \pts{10}
\[
A=\begin{bmatrix}
1&2&-1\\
3&10&0\\
-2&4&13
\end{bmatrix}.
\]
\workarea{1.7in}
\solution{Elimination uses $m_{21}=3$, $m_{31}=-2$, and then $m_{32}=2$. Hence
\[
\boxed{L=\begin{bmatrix}1&0&0\\3&1&0\\-2&2&1\end{bmatrix}},
\qquad
\boxed{U=\begin{bmatrix}1&2&-1\\0&4&3\\0&0&5\end{bmatrix}}.
\]}

\item Let
\[
A=\begin{bmatrix}
1&1&0\\
0&2&1\\
0&0&3
\end{bmatrix}.
\]
Define $B=3A$ and $C=B+I_3$. Compute the following determinants. \pts{10}
\begin{enumerate}[label=(\alph*)]
  \item $\det(-A)$ \hfill[4 pt]
  \item $\det B$ \hfill[3 pts]
  \item $\det C$ \hfill[3 pts]
\end{enumerate}
\workarea{1.6in}
\solution{Since $A$ is upper triangular, $\det A=1(2)(3)=6$.
\begin{enumerate}[label=(\alph*)]
  \item Since $A$ is $3\times3$, $\det(-A)=(-1)^3\det A=\boxed{-6}$.
  \item $\det B=\det(3A)=3^3\det A=27(6)=\boxed{162}$.
  \item $C=3A+I_3=\begin{bmatrix}4&3&0\\0&7&3\\0&0&10\end{bmatrix}$, so $\det C=4(7)(10)=\boxed{280}$.
\end{enumerate}}

\item Calculate the determinant of the following sparse $4\times4$ matrix. \pts{10}
\[
D=\begin{bmatrix}
1&0&2&0\\
0&3&0&0\\
4&0&-1&0\\
0&2&0&5
\end{bmatrix}.
\]
\workarea{2.7in}
\solution{Expand along column 4, and then along the second row of the remaining $3\times3$ determinant:
\[
\det D
=5\begin{vmatrix}1&0&2\\0&3&0\\4&0&-1\end{vmatrix}
=5(3)\begin{vmatrix}1&2\\4&-1\end{vmatrix}
=15(-1-8)=\boxed{-135}.
\]}

\newpage
\item Consider the system with parameters $a,b\in\R$:
\[
\begin{cases}
x+y+z=2,\\
2x+3y+4z=5,\\
3x+4y+(a+5)z=b.
\end{cases}
\]
Answer all parts. \pts{10}
\begin{enumerate}[label=(\alph*)]
  \item Use elimination to reduce the augmented matrix to row echelon form. \hfill[4 pts]
  \item Classify the system for all values of $a$ and $b$: state when it has one solution, no solution, or infinitely many solutions. \hfill[3 pts]
  \item Solve the system when $a=2$ and $b=9$. \hfill[3 pts]
\end{enumerate}
\workarea{3.0in}
\solution{
\begin{enumerate}[label=(\alph*)]
  \item Apply $R_2\leftarrow R_2-2R_1$, $R_3\leftarrow R_3-3R_1$, and then $R_3\leftarrow R_3-R_2$:
  \[
  \left[\begin{array}{ccc|c}
  1&1&1&2\\
  2&3&4&5\\
  3&4&a+5&b
  \end{array}\right]
  \sim
  \boxed{\left[\begin{array}{ccc|c}
  1&1&1&2\\
  0&1&2&1\\
  0&0&a&b-7
  \end{array}\right]}.
  \]
  \item If $a\ne0$, the system has a pivot in every variable column, so it has $\boxed{\text{one solution}}$. If $a=0$ and $b\ne7$, the last row represents $0=b-7$, so there is $\boxed{\text{no solution}}$. If $a=0$ and $b=7$, the last row is zero and $z$ is free, so there are $\boxed{\text{infinitely many solutions}}$.
  \item When $a=2$ and $b=9$, the last equation gives $2z=2$, so $z=1$. Then $y+2z=1$ gives $y=-1$, and $x+y+z=2$ gives $x=2$. Thus $\boxed{(x,y,z)=(2,-1,1)}$.
\end{enumerate}}
\end{enumerate}

\ifshowsolutions
\newpage
\section*{Compact answer key}
\begin{description}[leftmargin=2.5em]
  \item[MC:] B, B, B, A, B, A, D, A, B, D.
  \item[Brief:] rank $3$, free columns $3,5$; $\begin{bmatrix}1&2\\3&-1\end{bmatrix}$; $X=I+A^{-1}B$; $-3/2$; $AB$ is $4\times2$.
  \item[Exercises:]
  \begin{enumerate}[label=\arabic*.,leftmargin=2em,itemsep=0.15em,topsep=0.25em]
    \item $(x,y,z)=\boxed{(2,-1,3)}$.
    \item $A^2+AB+I=\boxed{\begin{bmatrix}3&3\\4&10\end{bmatrix}}$.
    \item $\displaystyle L=\boxed{\begin{bmatrix}1&0&0\\3&1&0\\-2&2&1\end{bmatrix}}$ and
    $\displaystyle U=\boxed{\begin{bmatrix}1&2&-1\\0&4&3\\0&0&5\end{bmatrix}}$.
    \item $\det(-A)=\boxed{-6}$, $\det B=\boxed{162}$, and $\det C=\boxed{280}$.
    \item $\det D=\boxed{-135}$.
    \item If $a\ne0$, there is one solution. If $a=0$ and $b\ne7$, there is no solution. If $a=0$ and $b=7$, there are infinitely many solutions. For $a=2$ and $b=9$, $(x,y,z)=\boxed{(2,-1,1)}$.
  \end{enumerate}
\end{description}
\fi
\newpage
Use the space below for scratch work or additional space for any of the questions above. 
\newpage
Additional scratch paper. You can detach this page for your convenience. \textbf{Please write your name and student ID on the top right of the page.}

\end{document}
